% =====================================================================
%  A Quantum Transformer World Model for Model-Based Reinforcement
%  Learning
%  V. Liutarevich
%
%  Draft v3 --- restructured 12 October 2026.
%
%  WHAT CHANGED FROM v2, AND WHY
%  -----------------------------
%  v2 led with a closed-loop result: under partial observation the
%  quantum agent led the classical control by 178.3 return points and was
%  ahead in three of three seeds. That claim is withdrawn. At seven
%  completed seeds the difference is +58.3 with a 95% interval of
%  [-111.9, +212.5] and an exact permutation p of 0.500. The dispersion
%  grew faster than the effect (sd 117 at three seeds, 237 at seven), so
%  detecting a difference of this size at 80% power would require roughly
%  130 seeds at 27 h each. CartPole returns cannot carry this paper.
%
%  The paper therefore now rests on:
%    (1) the architecture, which is the actual contribution;
%    (2) offline predictive performance, where a seed costs ~30 min
%        rather than 27 h, so eight seeds and a context ablation are
%        affordable;
%    (3) closed-loop viability reported honestly as parity;
%    (4) hardware execution.
%
%  Two corrections of record are marked in the text: the v2 imagination
%  table was computed before the symlog fix and is superseded, and the
%  v2 full-observation table contained two runs destroyed by a tag
%  collision.
%
%  PLACEHOLDERS
%  ------------
%     \pend{...}        inline, for a missing number or citation
%     \pendpara{...}    a paragraph-sized gap
%     pendingbox env    a whole section awaiting an experiment
%
%  \draftfalse hides every one of them, which is also how to see what
%  the paper looks like with the gaps removed. Appendix A lists them.
% =====================================================================
\documentclass[11pt]{article}

\usepackage[margin=1in]{geometry}
\usepackage[T1]{fontenc}
\usepackage[utf8]{inputenc}
\usepackage{amsmath,amssymb}
\usepackage{booktabs}
\usepackage{graphicx}
\usepackage{xcolor}
\usepackage{caption}
\usepackage[colorlinks=true,linkcolor=NavyBlue,citecolor=NavyBlue,
            urlcolor=NavyBlue]{hyperref}
\usepackage[numbers,sort&compress]{natbib}

\definecolor{NavyBlue}{HTML}{1F3864}
\definecolor{PendBG}{HTML}{FFF4E5}
\definecolor{PendFG}{HTML}{8A4B00}

% ---- placeholder machinery -------------------------------------------
\newif\ifdraft
\drafttrue                    % \draftfalse to hide every placeholder

\newcommand{\pend}[1]{%
  \ifdraft\textcolor{PendFG}{\textbf{[TODO:} \textit{#1}\textbf{]}}\fi}

\newcommand{\pendpara}[1]{%
  \ifdraft\par\vspace{2pt}\noindent
  \fcolorbox{PendFG}{PendBG}{%
    \begin{minipage}{0.97\linewidth}\small\textcolor{PendFG}{%
      \textbf{TO DO.} #1}\end{minipage}}\par\vspace{4pt}\fi}

% A named savebox rather than LaTeX's internal \@tempboxa, so the
% definition needs no \makeatletter and cannot clash with anything else.
\newsavebox{\pendbox}
\newenvironment{pendingbox}[1]
  {\par\vspace{3pt}\noindent
   \begin{lrbox}{\pendbox}\begin{minipage}{0.93\linewidth}\small
   \textcolor{PendFG}{\textbf{SECTION PENDING --- #1}}\par\vspace{3pt}
   \color{PendFG}}
  {\end{minipage}\end{lrbox}%
   \ifdraft\fcolorbox{PendFG}{PendBG}{\usebox{\pendbox}}\fi
   \par\vspace{4pt}}

\title{\bfseries A Quantum Transformer World Model for\\
       Model-Based Reinforcement Learning}
\author{V. Liutarevich}
\date{Draft v3 --- \today}

\begin{document}
\maketitle

\begin{abstract}
Quantum reinforcement learning has so far placed parameterised quantum
circuits in the policy or the value function. We place one in the
\emph{world model}, which we argue is the better-motivated target: a
world model predicts dynamics, and dynamics is what quantum systems
natively express.

We present, to our knowledge, the first \emph{quantum transformer world
model}: a transformer over a window of recent latent states, with a
parameterised quantum circuit as the token-mixing block of one encoder
layer, serving as the dynamics prior of a recurrent state-space model
and trained end to end inside a model-based agent. Because latent
imagination rolls the model forward on the prior alone, with no access
to observations, \textbf{the policy is trained entirely on trajectories
generated through a quantum circuit.}

Our primary evidence is offline: on held-out latent trajectories the
quantum transformer matches its capacity-matched classical twin at
one-step prediction and stays within the pre-registered factor of two at
ten-step open-loop rollout, over eight seeds.
\pend{replace with the measured figures from the width-matched context
sweep; see Section~\ref{sec:offline}}
Closed-loop, the architecture trains a competent policy and returns are
at parity; we report the dispersion that makes any stronger closed-loop
claim unattainable at this scale, and quantify it rather than assert it.
Measured directly, the quantum prior is the \emph{better} open-loop
dynamics model under partial observation, by $26\%$. Both circuits were
executed on IonQ Forte-1.

We make no claim of quantum advantage, and explain why none is available
at simulable scale.
\end{abstract}

% =====================================================================
\section{Introduction}
% =====================================================================
Reinforcement learning with quantum components has, to date, been almost
entirely model-free. Parameterised quantum circuits (PQCs) have been
used as policies and as value-function approximators, inheriting the
sample inefficiency of the model-free setting while adding the cost of
circuit evaluation.

Model-based reinforcement learning offers a different and, we argue,
better insertion point. A model-based agent learns a \emph{world
model} --- a predictive model of environment dynamics --- and optimises
its policy against that model rather than against the environment. The
world model is the component whose job is to represent dynamics, which
is what a quantum system does naturally.

This paper places a quantum transformer inside the world model of a
model-based agent and trains the whole system end to end. The
construction has two parts that are usually studied separately. From the
quantum sequence-modelling literature we take the idea of a transformer
in which one encoder layer's token-mixing block is a
PQC~\cite{qasa}. From the model-based RL literature we take the
recurrent state-space model (RSSM) of the DreamerV3
family~\cite{hafner2023}. We join them by making the quantum
transformer the RSSM's \textbf{prior}: the distribution over the next
stochastic latent given the recurrent state.

That join is the contribution. In latent imagination the agent rolls the
model forward using the prior alone, without observations, and trains
its policy on the resulting trajectories. Replacing the prior therefore
means the agent's imagined experience --- all of the data its policy
ever sees during optimisation --- is generated through a quantum
circuit. No prior work we are aware of puts a quantum model in that
position.

\paragraph{Why the evidence is where it is.}
A reader may expect the headline result of a paper like this to be a
closed-loop learning curve showing the quantum agent ahead. We do not
offer one, and the reason is quantitative rather than rhetorical. Over
seven completed seeds the paired closed-loop difference on partial-observation
CartPole is $+58.3$ return points with a standard deviation of $237$;
detecting an effect of that size at $80\%$ power needs on the order of
$130$ seeds, and a seed costs $27$ hours. Offline prediction error on
the same latent dynamics has a seed cost of about half an hour. We
therefore place the quantitative claims where the measurements can
actually be made, and report the closed-loop behaviour as what it
is: evidence that the architecture works, at parity, with the power
limitation stated.

\paragraph{Contributions.}
\begin{enumerate}\itemsep2pt
\item A quantum transformer dynamics prior for recurrent state-space
      models, and the first model-based agent whose world model is a
      quantum transformer, trained end to end
      (Section~\ref{sec:arch}).
\item An account of what the integration requires: a bounded history
      window with correct reset semantics across sequence, rollout and
      episode boundaries; mapping circuit readouts to categorical latent
      logits; maintaining gradient flow through multi-step imagination;
      and the interaction with KL balancing
      (Sections~\ref{sec:arch}--\ref{sec:impl}).
\item A verification harness targeting failure modes that are silent
      rather than loud, with the specific checks enumerated
      (Section~\ref{sec:impl}).
\item Offline predictive performance over eight seeds, with a
      width-matched ablation over attention context length
      (Section~\ref{sec:offline}).
\item A direct measurement of world-model quality showing that agent
      performance is only weakly coupled to it, and that the direction of
      the quality difference depends on observability
      (Section~\ref{sec:imagination}).
\item Execution of both circuits on trapped-ion hardware, with the
      measurement convention settled empirically rather than assumed
      (Section~\ref{sec:hardware}).
\end{enumerate}

We make no general claim of advantage. At the scale accessible to exact
simulation the circuit computes a truncated Fourier series whose
frequency support is fixed by the encoding~\cite{schuld2021}, and a
classical surrogate spanning the same spectrum reaches the same error
with far fewer parameters~\cite{papera}. Parity is the honest
expectation, and parity is what we find.

% =====================================================================
\section{Background: recurrent state-space models}
% =====================================================================
A model-based agent of the DreamerV3 family~\cite{hafner2023} maintains
a latent state factorised into a deterministic recurrent component
$h_t$ and a stochastic component $z_t$ represented as categorical
variables. Three distributions define the world model: a recurrent
update $h_t = f(h_{t-1}, z_{t-1}, a_{t-1})$; a \emph{posterior}
$q(z_t \mid h_t, o_t)$, which sees the observation; and a \emph{prior}
$p(z_t \mid h_t)$, which does not. Training reconstructs observations
and rewards while keeping the prior close to the posterior under a
balanced KL term.

\paragraph{Imagination.} Policy optimisation does not interact with the
environment. The agent samples a starting latent from replayed
experience and rolls the world model forward for a fixed horizon using
\emph{the prior only}, since no future observations exist. Actor and
critic are trained on these imagined trajectories via
$\lambda$-returns.

The prior is therefore the generative core of the agent: it is what the
agent dreams with, and the most consequential place to introduce a
quantum component.

\paragraph{Why a transformer prior is not redundant.}
The RSSM already has a recurrent pathway, so one might ask what
attention adds. The answer is that $h_t$ is a \emph{summary} produced by
a GRU trained under the world-model loss, and the prior sees only its
current value. A prior that attends over $h_{t-K+1}, \dots, h_t$ can
condition on the recent trajectory of the recurrent state rather than
its instantaneous value, which matters precisely when the dynamics are
history-dependent. This is the same argument that motivates
TransDreamer~\cite{transdreamer}, applied to the prior alone rather
than to the whole recurrent backbone --- a deliberately smaller
intervention, so that the comparison against the baseline isolates one
module.

% =====================================================================
\section{Related work}
% =====================================================================
\paragraph{Quantum reinforcement learning.} Existing work places
parameterised quantum circuits in the policy or value
function~\cite{chen2020,skolik2023}. To our knowledge no prior work
places a quantum circuit in the dynamics model of a model-based agent.
\pend{the systematic novelty search has been completed; record the
search strings, databases and dates in Appendix~\ref{app:novelty}}

\paragraph{World models for quantum problems.} The converse direction
exists: model-based reinforcement learning applied \emph{to} quantum
problems, for instance DreamQAS~\cite{dreamqas}, which learns a
classical world model over circuit-construction trajectories. That work
uses a classical model to design quantum circuits; we use a quantum
circuit inside a model.

\paragraph{Quantum attention.} QASA~\cite{qasa} introduces the
construction our prior adapts: a transformer encoder in which one
layer's token mixing is performed by a PQC. Their reported gain is, in
the revised version, attributed to the low-rank bottleneck the circuit
imposes rather than to quantumness as such --- a distinction we take
seriously, and which motivates the head-width ablation of
Section~\ref{sec:ablation}.

\paragraph{Transformers as world models.} TransDreamer~\cite{transdreamer}
replaces Dreamer's recurrent backbone with a transformer state-space
model. We leave the backbone intact and make only the prior a
transformer, which keeps every other component byte-identical to the
validated baseline.

\paragraph{Expressivity and benchmarking.} Data re-uploading circuits of
the kind we use compute truncated Fourier series in their inputs, with
frequencies fixed by the encoding gates and only the coefficients
trainable~\cite{schuld2021}. This bounds what the circuit can represent
and predicts that a classical model spanning the same frequency support
should match it; we confirm this directly in the companion
paper~\cite{papera}. Bowles et al.~\cite{bowles2024} find that quantum
models at simulable scale generally match rather than exceed classical
baselines of comparable capacity, and report that removing entanglement
often leaves performance unchanged --- which informs both our
expectations and our choice of controls.

% =====================================================================
\section{Architecture}\label{sec:arch}
% =====================================================================
\subsection{The quantum transformer prior}
In the baseline RSSM the prior is a multilayer perceptron mapping the
deterministic recurrent state to logits over the categorical stochastic
latent,
\[
  \text{logits} = \mathrm{MLP}(h_t) \in \mathbb{R}^{S \times C},
  \qquad p(z_t \mid h_t) = \mathrm{Categorical}(\text{logits}),
\]
for $S$ stochastic variables with $C$ classes each. We replace this map
with a quantum transformer over a bounded window of recent recurrent
states. Writing $H_t = (h_{t-K+1}, \dots, h_t)$ for the window, the
prior computes
\[
  \text{logits} = W_{\text{out}}\,
    \mathrm{Enc}_{\text{quantum}}\!\big(\mathrm{LN}(W_{\text{in}} H_t)
      + \mathrm{PE}\big)_{[-1]},
\]
where the final token's representation is read out, $\mathrm{PE}$ is
sinusoidal positional encoding, and $\mathrm{Enc}_{\text{quantum}}$ is a
two-block transformer encoder in which block $1$ carries the PQC. The
interface $h_t \mapsto \text{logits}$ is preserved exactly, so every
other component --- recurrent update, posterior, decoders, actor,
critic, KL balancing, the imagination loop --- is unchanged.

Each encoder block applies causally-masked multi-head self-attention,
then (in the quantum block) the circuit as a residual token-mixing
sublayer, then a position-wise feed-forward network, each with layer
normalisation:
\begin{align*}
  u &= \mathrm{LN}_1\big(x + \mathrm{MHA}(x, x, x)\big), \\
  v &= \mathrm{LN}_2\big(u + \mathrm{QuantumToken}(u)\big), \\
  y &= \mathrm{LN}_3\big(v + \mathrm{FFN}(v)\big).
\end{align*}
$\mathrm{QuantumToken}$ acts on each sequence position independently:
project $d_{\text{model}} \to n$ rotation angles via a linear map,
$\tanh$ and scaling by $\pi$; evaluate the circuit; project the
$n$ readouts back to $d_{\text{model}}$.

\begin{table}[h]\centering\small
\caption{The dynamics prior. The window length $K$ is the attention
context; $K = 4$ is used for all closed-loop runs.}
\label{tab:arch}
\begin{tabular}{lll}
\toprule
& Quantum transformer prior & Classical twin \\
\midrule
Attention context $K$        & 4                 & 4 \\
$d_{\text{model}}$           & 64                & 64 \\
Attention heads              & 4                 & 4 \\
Encoder blocks               & 2                 & 2 \\
Token-mixing block index     & 1                 & 1 \\
Token-mixing map             & PQC, 6 qubits     & $\tanh$ MLP, width 6 \\
Data qubits $+$ ancilla      & $6 + 1$           & --- \\
Circuit layers $L$           & 4                 & --- \\
Encoding                     & angle, re-uploaded & --- \\
Entangler                    & brickwork         & --- \\
Readout                      & local Pauli-$Z$   & --- \\
\midrule
Trainable circuit parameters & $L \times (n+1) = 28$ & 0 \\
Total prior parameters       & 84{,}770          & \pend{count} \\
\bottomrule
\end{tabular}
\end{table}

\subsection{The window, and why it is the fragile part}
The prior is called once per RSSM step by code that knows nothing about
attention, so the window is maintained as mutable state on the module
and advanced on each call. This is the component most likely to be
wrong in a way that produces a plausible result rather than an error. A
window that is not cleared at an episode boundary lets one episode
attend over the tail of the previous one; a window that is not cleared
between imagination rollouts lets imagined futures attend over replayed
history. Either leak would inflate apparent performance silently.

We therefore install the reset logic by wrapping the three entry points
at which a boundary occurs --- sequence start in \texttt{observe},
per-step reset in \texttt{obs\_step}, and rollout start in
\texttt{imagine} --- on the agent instance, leaving the baseline
implementation untouched. The semantics are asserted by test, not
assumed; see Section~\ref{sec:impl}.

\subsection{The bottleneck, and an honest accounting}
The circuit is $28$ of the prior's $84{,}770$ parameters: $0.03\%$. We
state this because it bounds what the architecture can be claimed to
demonstrate.

The defensible claim is one of \emph{information bottleneck} rather than
parameter share: in the quantum block, every token's representation is
routed through $n = 6$ real numbers produced by a quantum circuit, and
the surrounding layers cannot route around it, because the sublayer is
on the residual path of the only block that mixes tokens that way.
Whether the circuit or the surrounding width carries the predictive work
is an empirical question, and we answer it in
Section~\ref{sec:ablation} rather than assuming the favourable answer.

For comparison, the earlier non-attentional prior --- retained as the
\texttt{mlp} ablation arm, and the configuration used for all seven
partial-observation seeds in Section~\ref{sec:closedloop} --- had the
following breakdown.

\begin{table}[h]\centering\small
\caption{Parameter breakdown of the non-attentional (\texttt{mlp})
prior, retained as an ablation arm.}
\begin{tabular}{lr}
\toprule
Component & Parameters \\
\midrule
Down-projection ($256 \rightarrow 6$) & 1{,}542 \\
\textbf{Quantum circuit} & \textbf{28} \\
Expansion head ($6 \rightarrow 64 \rightarrow 256$) & 17{,}088 \\
\midrule
Total, quantum prior & 18{,}658 \\
Classical control, same positions & 18{,}672 \\
Original DreamerV3 prior (replaced by both) & 131{,}840 \\
\bottomrule
\end{tabular}
\end{table}

\subsection{Why the prior}
Three positions could host a quantum circuit: the policy, the value
function, or the world model. The first two are what the literature has
used. We take the world model because it is the component that models
dynamics; because the policy trains entirely on imagined rollouts, so
the prior has the greatest leverage over final performance; and because
it is the smallest intervention with the largest reach.

\subsection{The classical control}
To isolate the quantum substrate from the architectural change
accompanying it, the control is the identical transformer with the PQC
sublayer replaced by a $\tanh$ MLP of the same bottleneck width,
occupying the same position in the same block. Everything else ---
attention, positional encoding, layer norms, feed-forward network,
window semantics, context length --- is shared code. The arms are
required to match within $15\%$ on parameter count, and the test suite
fails if they do not.

This follows the capacity-matched protocol of~\cite{qasa}. Both arms
have a narrower bottleneck at that position than DreamerV3's original
prior, which is a property of the comparison to keep in view: we are
comparing two bottlenecked priors to each other, not either to a
full-width baseline.

% =====================================================================
\section{Implementation and verification}\label{sec:impl}
% =====================================================================
Four aspects of the integration proved non-obvious.

\paragraph{Optimiser reconstruction.} The prior is installed by
substituting a module within an already-constructed agent. An optimiser
built before the substitution holds references to the \emph{replaced}
parameters; the circuit parameters are then silently untracked, and
training proceeds without error while the quantum component never
updates. The optimiser must be rebuilt after substitution. This failure
is silent and produces plausible learning curves.

\paragraph{Gradient flow through imagination.} The actor is trained on
trajectories produced by rolling the prior forward for $H$ steps, so
gradients reach the circuit through a composition of $H$ applications of
the same circuit. We instrument and report gradient magnitudes at the
circuit parameters separately from the surrounding classical ones.

\paragraph{Window reset semantics.} Discussed in
Section~\ref{sec:arch}. The non-obvious cases are a batch-size change
between the observation and imagination phases, which must re-seed
rather than reuse the window, and a partial reset in which only some
rows of the batch have hit an episode boundary.

\paragraph{KL balancing.} The world-model loss regularises the prior
toward the posterior with asymmetric weighting. A prior with a
materially different parameterisation changes the effective strength of
this term.
\pend{report KL trajectories for both arms}

\paragraph{Verification harness.} Nine checks must pass before any
training run, for both arms:

\begin{enumerate}\itemsep1pt
\item Output shape and finiteness of the logits.
\item The window is bounded at $K$ after more than $K$ calls.
\item History changes the prediction --- otherwise the attention is
      decorative.
\item \texttt{reset()} fully clears, so two fresh sequences agree
      exactly.
\item Residual nondeterminism in training mode is attributable to
      attention dropout and not to retained state: the module is
      verified stochastic under \texttt{train()} and deterministic under
      \texttt{eval()}.
\item \texttt{reset\_rows()} is selective and complete --- flagged rows
      cleared, unflagged rows bit-identical.
\item A batch-size change re-seeds the window rather than reusing it.
\item Gradients reach the circuit parameters specifically, with non-zero
      magnitude, not merely the surrounding layers.
\item The two arms are capacity-matched within tolerance.
\end{enumerate}

Checks 3 and 4 must run under \texttt{eval()}: with dropout active, two
identical calls differ regardless of whether the window was cleared, so
check 3 would pass on noise and check 4 would fail on a correct
implementation. We record this because we made that error first, and it
is the kind of test bug that produces false confidence in both
directions at once.

The motivation for the harness is that the failure modes here are
silent: a missing gradient path, an untracked parameter group or a
leaking window produces runs that complete normally and report
plausible numbers.

\paragraph{Run naming.} Every configuration setting appears
unconditionally in the run directory name, and training refuses to write
into a directory containing a completed run. This is a response to an
incident rather than a precaution: an earlier scheme appended only
non-default values, so a head-width ablation at the default width
produced the same directory name as a headline run at a different step
budget and overwrote two completed $80{,}000$-step runs. Both runs
finished normally and left plausible result files behind, and the loss
was invisible until the numbers were checked against their file
timestamps.
\pend{note in Section~\ref{sec:closedloop} which full-observation seeds
are from the re-run}

% =====================================================================
\section{Offline predictive performance}\label{sec:offline}
% =====================================================================
The prior's job is to predict latent dynamics. We measure that directly
and in isolation, on offline trajectories, before asking what it does
inside an agent. This is where the statistical power is: a seed costs
roughly half an hour rather than $27$ hours, so eight seeds and a
context ablation are affordable.

\paragraph{Setup.} Latent trajectories are collected from a trained
DreamerV3 agent on CartPole under a mixed policy ($70\%$ greedy,
$30\%$ $\epsilon$-random) to give the state coverage a dynamics model
needs; a purely greedy policy yields a narrow on-policy manifold that
flatters every model equally. Dreamer's $[h_t \mid z_t]$ features are
compressed to six dimensions by an autoencoder, one feature per qubit
under angle encoding, with held-out explained variance of approximately
$91\%$. Splits are by episode, not by window, since windowing across a
split boundary leaks future frames into training.

Two metrics: one-step prediction error, and $k$-step \emph{open-loop}
rollout error in which the model re-consumes its own predictions. The
second is the regime a world model is actually used in; teacher-forced
error looks far better and is not what the imagination loop does. The
pre-registered gate is to match the capacity-matched classical twin at
one step and stay within $2\times$ at ten steps.

\begin{pendingbox}{headline offline table, eight seeds}
Awaiting the width-matched sweep. Required here: one-step and $k$-step
open-loop error for the quantum transformer and its classical twin at
$k \in \{1, 5, 10, 20\}$, over eight seeds, with geometric-mean ratios,
bootstrap intervals and exact permutation $p$-values on the paired
log-ratios.

The three-seed measurement at $K = 2$ gives a geometric-mean
quantum/classical ratio of $0.983$ at one step and $0.983$ at ten
steps --- parity --- with per-seed ratios spanning $0.81$ to $1.16$.
That spread across three seeds is the reason this table needs eight.
\end{pendingbox}

\subsection{Attention context length}
The architectural claim is that attending over recent recurrent states
helps when the dynamics are history-dependent. That predicts a trend:
the quantum/classical ratio, or at least the absolute error, should
improve with $K$ and then saturate.

\begin{pendingbox}{context ablation, $K \in \{2,4,8,16\}$ at fixed width}
\textbf{Correction of record.} The context sweep initially present on
disk is unusable: $K = 2$ was run at $d_{\text{model}} = 32$, $K = 4$ at
$8$, and $K = 8$ at $6$. Context is confounded with width across those
three points, so they cannot be compared and no trend can be read from
them. The sweep is being re-run at a single fixed width
($d_{\text{model}} = 32$, two blocks) over eight seeds.

This is the same class of error as the run-naming incident of
Section~\ref{sec:impl}: a setting varied without being recorded as
varied. The run-naming fix prevents silent overwriting but not silent
confounding, so the sweep driver now refuses to start when any
result-affecting setting absent from the directory name differs from its
default.
\end{pendingbox}

% =====================================================================
\section{Closed-loop behaviour}\label{sec:closedloop}
% =====================================================================
\subsection{Setup}
CartPole, with a reduced training configuration relative to DreamerV3's
defaults: batch $8$ against $16$, sequence length $16$ against $32$,
imagination horizon $10$ against $15$, and a gradient step every four
environment steps rather than every two. The reduction is forced by the
cost of circuit evaluation --- imagination performs \texttt{horizon}
sequential circuit calls over $\texttt{batch} \times
\texttt{seq\_len}$ states per gradient step --- and applies identically
to both arms. The solved threshold is a mean evaluation return of $475$.

Final returns are measured over $50$ evaluation episodes from the saved
checkpoint, rather than the $20$-episode figure recorded during
training. At $20$ episodes the standard error is large enough to move
individual seeds by up to $90$ points, which is comparable to the
differences being discussed.

\paragraph{Reference implementation.} Full DreamerV3, at its own default
settings, solves CartPole in all three seeds we ran, first crossing the
threshold at $10{,}000$, $15{,}000$ and $45{,}000$ environment steps.
Our reduced configuration therefore costs roughly a factor of two to
eight in sample efficiency, which is the context for neither arm below
reaching the gate at $80{,}000$ steps.

\paragraph{Attention context.} All closed-loop runs use $K = 4$. At
$K = 8$ the measured throughput is $0.4$ environment steps per second
against $2.8$ at $K = 4$, a factor of seven, which is the difference
between a $55$-hour run and an $8$-hour one.

\subsection{The architecture trains}
\begin{pendingbox}{quantum transformer closed-loop runs}
The quantum transformer prior has been verified to train
(all nine integration checks pass; a $2{,}500$-step run reaches a return
of $62.0$ against a random-policy return of approximately $22$) but the
full $80{,}000$-step runs have not yet been done. Five seeds per arm are
in progress at $K = 4$.

Required here: final returns per seed for both arms, learning curves
with interquartile bands, and the gradient-magnitude trace at the
circuit parameters.

\textbf{This section must not be written as a comparison.} See
Section~\ref{sec:power}: at the dispersion measured over seven seeds of
the non-attentional prior, five seeds cannot resolve a difference of the
size observed. The claim available from these runs is that the
architecture trains a competent policy and that returns are at parity.
\end{pendingbox}

\subsection{The non-attentional prior, for reference}
The seven-seed comparison below uses the \texttt{mlp} prior
(Section~\ref{sec:arch}), which is the configuration for which we have
the most seeds. It is reported because it is the dataset behind the
power analysis of Section~\ref{sec:power}, not as a result about the
paper's architecture.

\begin{table}[h]\centering\small
\caption{Final return at $80{,}000$ environment steps, $50$ evaluation
episodes, non-attentional prior. Partial observation: both velocity
components hidden, so the agent must infer velocity from history.}
\label{tab:po}
\begin{tabular}{lccc}
\toprule
Seed & Quantum & Classical & Difference \\
\midrule
0 & $369.7$ & $290.0$ & $+79.7$ \\
1 & $497.6$ & $242.2$ & $+255.4$ \\
2 & $386.8$ & $85.4$  & $+301.4$ \\
3 & $152.9$ & $352.7$ & $-199.8$ \\
4 & $328.7$ & $111.7$ & $+217.0$ \\
5 & $233.6$ & $155.0$ & $+78.6$ \\
6 & $70.9$  & $395.1$ & $-324.2$ \\
\midrule
Mean & $291.5$ & $233.2$ & $+58.3$ \\
\bottomrule
\end{tabular}
\end{table}

The paired difference is $+58.3$ with a $95\%$ bootstrap interval of
$[-111.9, +212.5]$ and an exact permutation $p$ of $0.500$; the quantum
arm is ahead in five of seven seeds. The interval contains zero. These
data do not establish a difference in either direction.

\begin{table}[h]\centering\small
\caption{Final return at $80{,}000$ steps, full observation,
non-attentional prior. Seeds $0$ and $1$ are from the re-run described
in Section~\ref{sec:impl}.}
\label{tab:full}
\begin{tabular}{lccc}
\toprule
Seed & Quantum & Classical & Difference \\
\midrule
0 & $85.7$  & $474.8$ & $-389.1$ \\
1 & $219.3$ & $302.2$ & $-82.9$ \\
2 & $235.5$ & $245.7$ & $-10.2$ \\
3 & $190.9$ & $483.0$ & $-292.1$ \\
4 & $166.9$ & $111.7$ & $+55.2$ \\
\midrule
Mean & $179.7$ & $323.5$ & $-143.8$ \\
\bottomrule
\end{tabular}
\end{table}

The paired difference is $-143.8$, interval $[-293.9, +1.4]$,
permutation $p = 0.188$, quantum ahead in one of five seeds. Again the
interval contains zero, though only just.

Full-state CartPole is Markovian --- the observation determines the next
state --- so a history-dependent prior has little to contribute there,
and the sign is in the direction the architecture would predict. We
decline to make more of it than that, because neither comparison is
resolved.

\pendpara{\textbf{Unexplained discrepancy, to resolve before
submission.} Full-observation seed $0$ returned $310.8$ in the original
run and $85.7$ in the re-run, at nominally identical configuration and
seed --- a factor of $3.6$. Seed $1$ moved from $247.7$ to $219.3$, a
factor of $1.1$. One plausible cause is that the two runs differed in
CPU thread count, which changes floating-point reduction order and so
the trajectory of an otherwise deterministic optimisation; the re-runs
were launched with \texttt{OMP\_NUM\_THREADS=1} and the originals may
not have been. This must be checked and either explained or reported: if
seed control does not survive a threading change, that is a
reproducibility fact about this pipeline which belongs in the paper,
and it bears on the companion paper's argument~\cite{papera}.}

\subsection{Why no stronger closed-loop claim is available}
\label{sec:power}
The partial-observation difference was $+212.2$ at three seeds,
$+122.1$ at six, and $+58.3$ at seven. The standard deviation of the
paired difference moved the other way: $117$, $182$, $237$. Dispersion
grew faster than the effect.

\begin{table}[h]\centering\small
\caption{The same comparison as seeds accumulate. $n_{80\%}$ is the
seeds required to detect the then-current effect at $80\%$ power and
$\alpha = 0.05$; $p_{\min}$ is the smallest two-sided $p$ attainable at
that sample size under an exact sign-flip permutation test.}
\begin{tabular}{lrrr}
\toprule
& 3 seeds & 6 seeds & 7 seeds \\
\midrule
Difference      & $+212.2$ & $+122.1$ & $+58.3$ \\
Ahead in        & $3/3$    & $5/6$    & $5/7$ \\
sd of difference& $117$    & $182$    & $237$ \\
Exact $p$       & $0.250$  & $0.156$  & $0.500$ \\
$p_{\min}$      & $0.250$  & $0.031$  & $0.016$ \\
$n_{80\%}$      & $3$      & $18$     & $\sim\!130$ \\
\bottomrule
\end{tabular}
\end{table}

At three seeds the smallest attainable two-sided $p$ is $0.250$, so that
configuration could not have reached significance whatever the data ---
a point worth making explicitly, because a three-seed study reporting
``ahead in all seeds'' invites a reading it cannot support. At $27$
hours per run, $130$ seeds is roughly five months of continuous compute
for one cell of one table.

This is not a peculiarity of our setup. The companion
paper~\cite{papera} measures seed dispersion across five sequence-modelling
tasks and finds the required sample size for a $10\%$ effect spanning
$6$ to $150$ depending on task, that under the paired null $29\%$ of
five-seed studies report a spurious $\ge 10\%$ gain, and that a
three-seed pilot cannot forecast its own required sample size: the same
pilot data is consistent with $n$ between $4$ and $170$. We therefore
treat the closed-loop comparison as a viability demonstration and put
the quantitative weight on Section~\ref{sec:offline}.

% =====================================================================
\section{Is the world model used?}\label{sec:imagination}
% =====================================================================
An agent can reach competence while effectively ignoring a poor world
model, learning a policy that works despite the predictions rather than
because of them. We therefore measure the prior directly: roll it
forward without observations, as imagination does, and compare the
decoded predictions against the observations actually received.

\paragraph{A measurement error worth recording.} Our first version of
this measurement compared the decoder's output against raw
observations. The DreamerV3 decoder is trained against
$\mathrm{symlog}(o)$, so that comparison applied an unmodelled
transformation to one side and biased the result against the model. The
figures reported in draft~v2 of this paper --- quantum $0.0371$ against
classical $0.0210$ under full observation, and $0.0070$ against
$0.0052$ under partial --- are superseded. All numbers below are in
symlog space, which is where the decoder is trained.

\begin{table}[h]\centering\small
\caption{Open-loop prediction error at horizon $10$, in symlog space,
mean over seeds $0$--$5$. Lower is better.}
\begin{tabular}{lccc}
\toprule
Condition & Classical prior & Quantum prior & Ratio \\
\midrule
Partial observation & $0.00264$ & $\mathbf{0.00195}$ & $0.74\times$ \\
Full observation    & \pend{re-measure} & \pend{re-measure} & \pend{} \\
\bottomrule
\end{tabular}
\end{table}

Under partial observation the quantum prior is the \emph{better}
open-loop dynamics model, by $26\%$ --- the condition in which the
dynamics are history-dependent and a predictive model has something to
do. This is the measurement that most directly supports the
architectural hypothesis, and it is also the measurement with the least
noise, because it is taken over many transitions rather than over
whole-episode returns.
\pend{add per-seed values and the paired permutation test; a mean over
six seeds without dispersion is exactly the sin Section~\ref{sec:power}
warns against}

Yet the closed-loop returns in the same condition are at parity
(Table~\ref{tab:po}). \textbf{Agent performance is only weakly coupled
to world-model accuracy.} The mechanism is visible in the architecture:
when acting, the agent uses the posterior, which sees real
observations; the prior only shapes the imagined trajectories used to
train the actor and critic. A better prior yields better imagined
training data, which is a second-order effect on returns.

We note that a persistence baseline --- predicting no change --- is
adversarially strong on a stabilised CartPole, since the agent's task is
precisely to hold the state still, and we therefore report absolute
error rather than a ratio to it.

% =====================================================================
\section{Design space}\label{sec:ablation}
% =====================================================================
\begin{pendingbox}{head-width and entangler ablations}
Both sweeps have run at $40{,}000$ steps and two seeds on the
non-attentional prior; the analysis is not written, and two seeds is
below what Section~\ref{sec:power} shows is interpretable. Required
here:
\begin{itemize}\itemsep1pt
\item \textbf{Bottleneck width.} The width surrounding the circuit,
      reducing the expansion head from $17{,}088$ parameters to
      $1{,}308$. Performance roughly flat as it shrinks would show the
      circuit carries the predictive work; a collapse would show the
      surrounding width does, and would require the bottleneck claim of
      Section~\ref{sec:arch} to be withdrawn. Given that
      QASA's~\cite{qasa} revised version attributes its own gain to the
      bottleneck rather than to quantumness, this is the ablation a
      referee will ask for first.
\item \textbf{Entanglement removed.} All two-qubit gates deleted,
      parameter count unchanged, so the comparison is otherwise exact.
      Bowles et al.~\cite{bowles2024} report that removing entanglement
      frequently leaves performance unchanged; an architecture paper
      should establish this of its own construction rather than leave it
      to a referee.
\item \textbf{Quantum block position.} Which encoder block carries the
      circuit. QASA's central claim is that position matters more than
      count, and we have adopted their finding without re-testing it in
      this setting.
\item \textbf{Qubit count, circuit depth, readout locality.} Not run,
      and deliberately out of scope: the companion paper~\cite{papera}
      measures the relevant trade-off directly, finding a
      $2{,}766\times$ difference in measurement cost between entangler
      topologies at ten qubits, which bounds what a width sweep could
      mean for a hardware implementation.
\end{itemize}
\end{pendingbox}

% =====================================================================
\section{Execution on quantum hardware}\label{sec:hardware}
% =====================================================================
Both circuits were executed on IonQ Forte-1 at eight data qubits (nine
wires), two layers, $504$ shots, using the parameter-shift rule: two
circuit evaluations per gradient, four jobs in total.

\begin{table}[h]\centering\small
\caption{Gradients measured on IonQ Forte-1 against exact simulation.}
\begin{tabular}{lccc}
\toprule
Entangler & Hardware & Exact & Difference \\
\midrule
Brickwork    & $-0.6230$ & $-0.5533$ & $-0.0697$ \\
Ring cascade & $+0.0734$ & $+0.0380$ & $+0.0354$ \\
\bottomrule
\end{tabular}
\end{table}

Both agree with exact simulation within the measurement error. Our
reconstruction of $\langle Z_0 \rangle$ from the returned probability
histograms reproduces the PennyLane plugin's own expectation values
exactly, as rational multiples of $1/504$, which settles the
bit-ordering convention empirically rather than by assumption --- worth
stating because an inverted convention would have produced a plausible
disagreement attributable to noise.

The measured RMS deviation from exact over the four evaluations is
$0.0594$, against $0.0345$ predicted by IonQ's own Forte-1 noise model
scaled to this shot count. We report this as indicative only: it is a
single RMS over four evaluations at one parameter point, carrying
roughly $35\%$ relative uncertainty, and the parameter point chosen has
a ring-cascade gradient about $2.5\times$ larger than the landscape
average. The hardware run validates that measurements agree with
simulation; it is not evidence about typical gradient magnitudes.

\paragraph{Why only four jobs.} Trapped-ion time is billed per job, and
these four cost \pend{amount} in total. Training on hardware is not
merely expensive but infeasible: at the parameter count and step count
of one $80{,}000$-step run, the parameter-shift rule requires on the
order of $10^{12}$ shots. Hardware execution here is inference-only and
is intended to establish that the circuits transfer, not that the
training does.

% =====================================================================
\section{Discussion}
% =====================================================================
\paragraph{What this establishes.} A quantum transformer can serve as
the generative dynamics model of a model-based agent, trained end to end
through multi-step imagination, with the resulting policy learning
entirely from trajectories generated through a quantum circuit. The
construction is stable, reaches competence, admits a capacity-matched
classical control differing in exactly one sublayer, and runs on
hardware. Measured as a dynamics model under partial observation it is
$26\%$ better than that control.

\paragraph{What it does not.} We make no claim of quantum advantage, and
the theory says none is available here: a data re-uploading circuit
computes a truncated Fourier series whose frequencies are fixed by the
encoding, and a classical surrogate spanning the same spectrum matches
it with eight parameters against the circuit's
thirty~\cite{schuld2021,papera}. Closed-loop returns are at parity, and
the sample sizes at which they could be shown otherwise are out of
reach at $27$ hours per seed.

\paragraph{Why the architecture matters anyway.} It is a prerequisite
for questions that cannot be asked without it: whether quantum dynamics
models help where the dynamics are \emph{themselves} quantum --- the
natural next target, where the Fourier argument no longer bounds the
advantage away --- whether they transfer to hardware, and how they
behave under finite measurement. The contribution is a working
instrument, honestly characterised, not a benchmark win.

% =====================================================================
\section{Limitations}
% =====================================================================
\begin{itemize}\itemsep2pt
\item A single architecture family, at small qubit count, simulated
      exactly during training. Hardware execution was inference-only.
\item Closed-loop evidence is a viability demonstration, not a
      comparison. Five seeds on full observation and seven on partial,
      against a requirement of roughly $14$ and $130$ respectively.
\item Evaluation on low-dimensional control only. Scaling to pixel
      observations, where world models are most valuable, is untested.
\item The offline context ablation is being re-run after a width
      confound (Section~\ref{sec:offline}); the superseded version
      should not be cited from any earlier draft.
\item One full-observation seed is not reproducible across runs by a
      factor of $3.6$, cause not yet established
      (Section~\ref{sec:closedloop}).
\item The quantum prior costs roughly $12\times$ the classical control
      in wall clock on the offline task ($\sim\!1{,}700$\,s against
      $\sim\!150$\,s per seed). \pend{give the closed-loop ratio from
      runs known to have run uninterrupted; the earlier figure of
      $5.4\times$ is contaminated by machine sleep}
\item All quantitative claims are on one environment. A second
      (Pendulum) exists in the codebase but its continuous-action paths
      are untested.
\end{itemize}

% =====================================================================
\section{Reproducibility}
% =====================================================================
Implementation, verification harness, configurations and run records at
\pend{repository URL and archival DOI}. The world model is a faithful
small-scale DreamerV3; the quantum transformer prior is a drop-in
substitution installed on an already-constructed agent, so the two arms
differ in exactly one sublayer and the baseline implementation is
unmodified. Run directories encode every configuration setting,
training refuses to write into a directory that already contains a
completed run, and the sweep driver refuses to start when a
result-affecting setting absent from the directory name differs from its
default.

We report the failures as well as the results --- the overwrite, the
symlog mismatch, the width confound, the irreproducible seed --- because
each was silent, each produced plausible numbers, and the reader has no
way to rule them out of a paper that does not mention them.

% =====================================================================
\appendix
\section{Work remaining}\label{app:remaining}
% =====================================================================
In priority order. Items marked \textbf{blocking} must be resolved
before submission.

\begin{enumerate}\itemsep3pt
\item \textbf{Blocking.} Quantum transformer closed-loop runs, five
      seeds per arm at $K = 4$, partial observation
      (Section~\ref{sec:closedloop}). Without these the paper's titular
      architecture has never been run closed-loop.
\item \textbf{Blocking.} Width-matched offline sweep, $K \in \{2,4,8\}$,
      eight seeds, both arms (Section~\ref{sec:offline}). The primary
      evidence.
\item \textbf{Blocking.} Confirm every partial-observation run in
      Table~\ref{tab:po} uses the same prior architecture. A window
      existed in which \texttt{prior\_arch} was live but the run name
      did not encode it.
\item \textbf{Blocking.} Per-seed values and paired test for the
      open-loop comparison of Section~\ref{sec:imagination}; re-measure
      the full-observation row post-symlog-fix.
\item Bottleneck-width and entanglement-removed ablations, at enough
      seeds to interpret (Section~\ref{sec:ablation}).
\item Resolve the seed-$0$ reproducibility discrepancy; test explicitly
      whether thread count changes trajectories.
\item Learning curves and circuit gradient-magnitude figure.
\item KL trajectories for both arms.
\item $K = 16$ as a fourth context point.
\item Record the novelty search: strings, databases, dates
      (Appendix~\ref{app:novelty}).
\item Repository URL, archival DOI, hardware cost figure, and full
      citation details for all references.
\end{enumerate}

% =====================================================================
\section{Novelty search}\label{app:novelty}
% =====================================================================
\begin{pendingbox}{novelty search record}
The search was run and cleared. Record here, so the first-of-its-kind
claim is auditable: databases queried, exact search strings, date of
each query, number of hits screened, and the specific near misses with
the reason each was excluded. The claim in the abstract is only as good
as this appendix.
\end{pendingbox}

% =====================================================================
\bibliographystyle{unsrtnat}
\begin{thebibliography}{9}
\bibitem{hafner2023} D.~Hafner, J.~Pasukonis, J.~Ba, T.~Lillicrap.
  \emph{Mastering Diverse Domains through World Models} (DreamerV3),
  arXiv:2301.04104. \pend{confirm venue and year of record}
\bibitem{transdreamer} C.~Chen, Y.-F. Wu, J.~Yoon, S.~Ahn.
  \emph{TransDreamer: Reinforcement Learning with Transformer World
  Models}, arXiv:2202.09481 (2022). \pend{verify author list}
\bibitem{chen2020} S.~Y.-C. Chen et al. \emph{Variational Quantum
  Circuits for Deep Reinforcement Learning}. \pend{full citation}
\bibitem{skolik2023} \emph{Variational Quantum Circuit Design for
  Quantum Reinforcement Learning on Continuous Environments},
  arXiv:2312.13798. \pend{verify authors}
\bibitem{dreamqas} \emph{DreamQAS: Learning a Decision-Useful World
  Model for VQE-Efficient Quantum Architecture Search},
  arXiv:2607.29491. \pend{verify authors}
\bibitem{qasa} \emph{Quantum Adaptive Self-Attention for Quantum
  Transformer Models}, arXiv:2504.05336. \pend{verify authors; cite the
  revised version, which attributes its reported gain to a low-rank
  bottleneck rather than to quantumness}
\bibitem{schuld2021} M.~Schuld, R.~Sweke, J.~J. Meyer. \emph{Effect of
  data encoding on the expressive power of variational quantum machine
  learning models}, Phys.\ Rev.\ A \textbf{103}, 032430 (2021).
\bibitem{bowles2024} J.~Bowles, S.~Ahmed, M.~Schuld. \emph{Better than
  classical? The subtle art of benchmarking quantum machine learning
  models}, arXiv:2403.07059 (2024).
\bibitem{papera} V.~Liutarevich. \emph{How Many Seeds? Dispersion,
  Resolving Power, and a Failed Replication in Quantum Sequence
  Modelling}. \pend{companion paper; cite as submitted or as preprint}
\end{thebibliography}

\end{document}
